\documentclass[12pt]{article}
\usepackage[margin=0.75in]{geometry}
\usepackage{fontspec}
\setmainfont{Latin Modern Roman}
\usepackage{microtype}
\usepackage{multicol}
\usepackage{fancyhdr}
\usepackage{titlesec}
\usepackage{enumitem}
\usepackage{xcolor}
\usepackage[hidelinks]{hyperref}
\setlength{\columnsep}{0.28in}
\setlength{\parindent}{1.1em}
\setlength{\parskip}{0.32em}
\setlength{\headheight}{15pt}
\titleformat{\section}{\large\bfseries\scshape}{}{0pt}{}
\titleformat{\subsection}{\normalsize\bfseries}{}{0pt}{}
\pagestyle{fancy}
\fancyhf{}
\fancyfoot[C]{\thepage}
\renewcommand{\headrulewidth}{0.4pt}
\newcommand{\focusbox}[1]{\par\vspace{0.4em}\noindent\begingroup\setlength{\fboxsep}{6pt}\colorbox{gray!12}{\begin{minipage}{\dimexpr\linewidth-2\fboxsep\relax}#1\end{minipage}}\endgroup\par\vspace{0.5em}}

\fancyhead[L]{Whately: Analogy}
\fancyhead[R]{Student Reading}
\begin{document}
\begin{center}
{\LARGE\bfseries Analogy and Relevant Similarity}\par\vspace{0.25em}
{\large\scshape Week 11: When Comparisons Prove Too Much}\par\vspace{0.5em}
{Adapted and modernized from Richard Whately, \textit{Elements of Logic}}\par\vspace{0.6em}
\rule{0.82\textwidth}{0.4pt}
\end{center}

\section*{Central Question}
When does a comparison actually prove something?

\section*{Vocabulary Preview}
\begin{multicols}{2}
\begin{itemize}[leftmargin=*, itemsep=0.18em]
\item \textbf{Analogy}: reasoning from a likeness between two cases.
\item \textbf{Source case}: the case used as the comparison.
\item \textbf{Target case}: the case the speaker wants us to judge.
\item \textbf{Relevant similarity}: a likeness that matters for the conclusion.
\item \textbf{Disanalogy}: an important difference between the compared cases.
\item \textbf{Overextension}: forcing a comparison beyond what it can prove.
\item \textbf{Decorative comparison}: a comparison that adds color but little proof.
\item \textbf{Analogy test}: claim, comparison, relevant likeness, disanalogy, judgment.
\item \textbf{Burden of fit}: the need to show that the comparison fits the disputed point.
\item \textbf{Narrow repair}: reducing the conclusion to what the analogy actually supports.
\end{itemize}
\end{multicols}

\section*{Introduction}
Some arguments compare one case to another: a school is like a household, a law is like a fence, a habit is like a muscle, a ruler is like a captain, a rumor is like a spark. Comparisons can help us reason. They can also make weak proof sound vivid.

Whately often reminds readers that probable reasoning depends on the real relation between cases. A likeness is not proof simply because it is memorable. Week 11 asks whether the similarity matters for the conclusion being drawn.

\focusbox{\textbf{Source note:} Adapted and modernized from Whately's treatment of comparison, probability, and practical reasoning. Classroom examples are original. This week connects fallacy work to student essays, where analogies can clarify a claim or distract from a missing proof.}

\begin{multicols}{2}
\section*{Reading}

\subsection*{1. Why Analogies Persuade}
An analogy borrows familiarity. If a hard case feels like an easier case, the mind moves quickly. That speed can be useful. A good comparison lets us see structure.

But speed can also hide a gap. The fact that two things are alike in one respect does not mean they are alike in the respect needed for the conclusion.

\subsection*{2. Source Case and Target Case}
Every analogy has two sides:
\begin{itemize}[leftmargin=*]
\item The \textbf{source case} is the familiar comparison.
\item The \textbf{target case} is the thing being argued about.
\end{itemize}
If someone says, ``A classroom is like a courtroom, so students should have a chance to answer charges before punishment,'' the courtroom is the source case; the classroom discipline case is the target case.

\subsection*{3. Relevant Similarity}
A similarity is relevant only if it matters for the conclusion.

In the classroom/courtroom example, both situations involve accusation, evidence, authority, and possible penalty. Those likenesses matter if the conclusion concerns fairness before punishment. The analogy would be weaker if the conclusion were, ``therefore classrooms should use judges' robes.'' Robes are not the relevant feature.

\subsection*{4. Disanalogy}
A disanalogy is an important difference. It does not automatically destroy the comparison, but it may limit the conclusion.

A school is like a workplace in some ways: deadlines, cooperation, responsibility. It is unlike a workplace in other ways: students are minors, grades are part of teaching, and the main purpose is formation rather than employment. A strong thinker names both sides.

\subsection*{5. Overextension}
An analogy overextends when it proves more than the comparison can bear.

``Studying is like athletic training, so daily practice matters'' may be strong. ``Studying is like athletic training, so every student should be ranked publicly after every exercise'' is a much larger claim. The comparison may support practice without supporting public ranking.

\subsection*{6. Decorative Comparison}
Some comparisons are vivid but not evidential. ``This policy is a train wreck'' may show dislike, urgency, or emotional force. By itself it does not prove which part of the policy fails, why it fails, or what repair would work.

Do not ban decorative comparisons. In literature and rhetoric they matter. But when a speaker offers a comparison as proof, ask whether it supplies a relevant likeness.

\subsection*{7. The Analogy Test}
Use this five-step test:
\begin{enumerate}[leftmargin=*,itemsep=0.12em]
\item State the conclusion.
\item Name the source case and the target case.
\item Identify the similarity the speaker relies on.
\item Name at least one relevant disanalogy.
\item Decide whether the analogy proves the claim, supports only a narrower claim, or merely decorates the argument.
\end{enumerate}

\subsection*{8. Literary Cases}
In \textit{Julius Caesar}, Brutus compares Caesar to a serpent's egg: harmless now, dangerous if allowed to hatch. The comparison is powerful because it makes future danger feel obvious. The logical question is whether Caesar is really like the egg in the relevant respect. In \textit{Macbeth}, images of sickness, clothing, animals, and traps can clarify a character's condition, but they may also tempt a reader to prove too much from a metaphor.

\subsection*{9. Repairing an Analogy}
There are two honest repairs:
\begin{itemize}[leftmargin=*]
\item Add evidence that the two cases are alike in the relevant respect.
\item Narrow the conclusion so it matches the likeness actually shown.
\end{itemize}
A disciplined thinker can admire a comparison and still ask what it proves.

\subsection*{10. What This Week Adds}
A complete Week 11 diagnosis should include:
\begin{itemize}[leftmargin=*,itemsep=0.12em]
\item the conclusion;
\item the source and target cases;
\item the relevant similarity;
\item the important disanalogy;
\item a judgment about whether the analogy proves, partly supports, or merely decorates the claim.
\end{itemize}

\section*{Reading Focus}
\begin{enumerate}[leftmargin=*]
\item What is the difference between the source case and the target case?
\item Why is not every similarity a relevant similarity?
\item What is a disanalogy, and why does it matter?
\item Give one original example of an analogy that supports a narrow conclusion but not a broad one.
\item List the five steps of the analogy test.
\item How can an analogy in an essay decorate a paragraph without proving the thesis?
\item In \textit{Julius Caesar}, why does the serpent's egg comparison need testing rather than automatic acceptance?
\end{enumerate}

\section*{Handout Summary}
Write 5--7 sentences explaining how to test an analogy. Your summary must include the words \textit{source}, \textit{target}, \textit{relevant}, and \textit{disanalogy}.
\end{multicols}
\end{document}
